\documentclass[aspectratio=169,11pt]{beamer}
\usetheme{PoliticsMenu}

\usepackage{booktabs}
\usepackage{pgfplots}
\pgfplotsset{compat=1.18}

\title[Politics on the Menu]{Politics on the Menu:}
\subtitle{Political Identity, Consumer Sorting,\\and Restaurant Engagement}
\author{Your Name}
\institute[Your University]{Department of Economics\\Your University}
\date{April 12, 2024}
\PMseminar{Economics Research Seminar}
\PMtopright{Economics Research Seminar\quad\textbar\quad Spring 2024}

\begin{document}

\begin{frame}
  \titlepage
\end{frame}

\section{Motivation}

\begin{frame}{Why Do Political Identities Matter in Everyday Choices?}
\begin{columns}[T,totalwidth=\textwidth]
  \begin{column}{0.45\textwidth}
    \begin{itemize}
      \item Political identity may influence not only what people think, but also where they choose to spend time and money.
      \item Restaurants are a salient social space where political identity can shape consumer behavior beyond explicit politics.
      \item We study restaurant choice, engagement, and consumer sorting across politically identified establishments.
    \end{itemize}
  \end{column}
  \begin{column}{0.50\textwidth}
    \begin{PMpanel}{Growth in Politically Identified Restaurant Reviews}
    \centering
    \begin{tikzpicture}
      \begin{axis}[
        width=\linewidth,height=4.15cm,
        xmin=2016,xmax=2022,ymin=0,ymax=210,
        axis lines=left,
        xtick={2016,...,2022},
        ytick={0,50,100,150,200},
        tick label style={font=\scriptsize,text=PMInk},
        legend style={draw=none,font=\scriptsize,at={(0.5,-0.20)},anchor=north,legend columns=2},
        every axis plot/.append style={thick}
      ]
      \addplot+[PMBlue,mark=*,mark size=1.8pt] coordinates {(2016,12)(2017,24)(2018,35)(2019,47)(2020,62)(2021,89)(2022,108)};
      \addlegendentry{Blue series}
      \addplot+[PMYellow,mark=*,mark size=1.8pt] coordinates {(2016,16)(2017,39)(2018,57)(2019,79)(2020,104)(2021,145)(2022,183)};
      \addlegendentry{Yellow series}
      \end{axis}
    \end{tikzpicture}
    \end{PMpanel}
  \end{column}
\end{columns}
\end{frame}

\section{Data}

\begin{frame}{Data and Measurement}
\begin{columns}[T,totalwidth=\textwidth]
  \begin{column}{0.52\textwidth}
    \begin{PMpanel}{OpenRice review archive}
      \begin{itemize}
        \item 2,269,004 district-resolved eligible reviews.
        \item 260,911 reviewers and 39,199 establishments.
        \item Research window: 2016--2025.
      \end{itemize}
    \end{PMpanel}
    \vspace{.28cm}
    \begin{PMtakeaway}
      The color labels are measurement devices. White does not mean politically neutral.
    \end{PMtakeaway}
  \end{column}
  \begin{column}{0.43\textwidth}
    \begin{PMpanel}{Three-color classification}
      \centering
      \begin{tabular}{lr}
        \toprule
        Group & Restaurants \\
        \midrule
        Yellow & 3,837 \\
        Blue & 2,469 \\
        White & 32,893 \\
        \bottomrule
      \end{tabular}
    \end{PMpanel}
  \end{column}
\end{columns}
\end{frame}

\section{Empirical Strategy}

\begin{frame}{Empirical Strategy and Main Results}
\begin{columns}[T,totalwidth=\textwidth]
  \begin{column}{0.48\textwidth}
    \begin{PMpanel}{Empirical Strategy}
      \begin{enumerate}
        \item Compare allocation across fixed restaurant color groups.
        \item Match restaurants around first qualifying Yellow reviews.
        \item Estimate conditional review choice across Yellow, Blue, and White restaurants.
      \end{enumerate}
    \end{PMpanel}
  \end{column}
  \begin{column}{0.48\textwidth}
    \begin{PMpanel}{Main Results}
      \centering\scriptsize
      \begin{tabular}{lcc}
        \toprule
        & Count & Rating \\
        \midrule
        Months +1 to +3  & 0.233*** & 0.983*** \\
                         & (0.031)  & (0.242)  \\
        Months +4 to +6  & 0.251*** & 1.289*** \\
                         & (0.036)  & (0.258)  \\
        Months +7 to +12 & 0.192*** & 1.031*** \\
                         & (0.042)  & (0.243)  \\
        \bottomrule
      \end{tabular}
      \vspace{.20cm}

      \raggedright\tiny Notes: illustrative layout using manuscript estimates.
    \end{PMpanel}
  \end{column}
\end{columns}
\end{frame}

\section{Results}

\begin{frame}{Conditional Review Choice}
\begin{PMpanel}{Main message}
Estimated availability makes early residual Yellow attraction model-sensitive; the reopening-period decline persists.
\end{PMpanel}

\vspace{.28cm}
\begin{columns}[T,totalwidth=\textwidth]
  \begin{column}{0.50\textwidth}
    \begin{itemize}
      \item Availability differs from proportional adjustment.
      \item Early attraction depends on supply normalization.
      \item The late negative Yellow index remains across specifications.
    \end{itemize}
  \end{column}
  \begin{column}{0.45\textwidth}
    \begin{PMpanel}{Model object}
      Conditional allocation across colors, holding district-quarter activity and recorded restaurant stock fixed.
    \end{PMpanel}
  \end{column}
\end{columns}
\end{frame}

\section{Conclusion}

\begin{frame}{Conclusion}
\begin{itemize}
  \item Market sorting, public identity expression, and business evaluation are related but not interchangeable.
  \item Local attention-associated gains can coexist with weaker later aggregate attraction.
  \item The evidence concerns recorded reviews and ratings, not political tastes, transactions, or welfare.
\end{itemize}
\vspace{.45cm}
\begin{PMtakeaway}
Use navy for structure, yellow for emphasis, and white space for the identification story.
\end{PMtakeaway}
\end{frame}

\end{document}
